\section{Use Cases}
\label{sec:usecases}

\subsection{SQL/PGQ with DuckPGQ}
\todo{Before: fork of DuckDB with patched libpg\_query grammar
(cwida/duckdb-pgq)~\cite{tenwolde2023duckpgq}. After: GrammarChange API, loadable
against upstream DuckDB. Lines of code before/after, maintenance cost (point
release merges), what SQL/PGQ syntax needed (fused operator tokens, label
expressions, quantifiers). Pitfall: wrapper rules without transform process.}

\subsection{SaneQL as a Loadable Dialect}
\todo{SaneQL~\cite{neumann2024saneql} as DialectExtension: \num{24} PEG rules,
\num{753} C++ lines, no core changes, all 22 TPC-H queries pass. Table from
saneql/PAPER\_NOTES.md: reference translator 5{,}028 lines (fair subset
grammar+lexer+semana 2{,}641) vs 753. Framing: syntax is a thin replaceable layer
when the semantic layer is shared. Report all caveats: \texttt{defun} not
implemented (closed operator keyword set), host SQL lexical decisions leak into
the dialect (reserved words, \texttt{||}, \texttt{a.desc()}), \texttt{let} as three
rewrites, API gaps (AddTerminalRuleOverride, undocumented parse-tree shape).}

\subsection{\todo{Third Use Case}}
\todo{Options: (a) agentic SQL extensions as parse-time rewrites via GrammarChange:
\texttt{INSPECT} (schema+stats+samples) and a \texttt{WITH OUTPUT} clause
(max\_rows, truncate, total count); (b) another DialectExtension user
(Spark dialect is private; GGSQL possible collaborator; pipe syntax~\cite{shute2024pipe}
as a small dialect would be a clean third point). Decide by what can be built and
measured in time.}
